\section{Personas}
\begin{personaenv}
  {Laura Ricci}
  {24}
  {Catania, Italia}
  {single}
  {Mobile App Developer}
  {Utente}
  {curiosa, collaborativa, precisa}
  {\begin{itemize}
    \item Segnalare bug in modo rapido e chiaro.
    \item Collaborare con altri sviluppatori in tempo reale.
    \item Ricevere notifiche sugli aggiornamenti.
  \end{itemize}}
  {\begin{itemize}
    \item Architetture distribuite.
    \item Strumenti collaborativi per il lavoro in team.
    \item Sviluppo mobile (Android/Kotlin).
  \end{itemize}}
  {Laura è una ragazza capace e disponibile che ha trovato lavoro come junior developer dopo aver conseguito
  la laurea in Informatica con il massimo dei voti. Specializzatasi successivamente nello sviluppo di app mobile,
  lavora in un piccolo team di ragazzi capaci e volenterosi di collaborare con la propria utenza per migliorare
  i propri prodotti.}
\end{personaenv}

\begin{personaenv}
  {Christian Ranavolo}
  {21}
  {Napoli, Italia}
  {fidanzato}
  {Junior Developer}
  {Utente}
  {arrendevole, riservato, scorbutico}
  {\begin{itemize}
    \item Monitorare bug delle app in sviluppo.
    \item Tenere traccia delle proprie attività.
    \item Filtrare rapidamente per priorità e stato.
  \end{itemize}}
  {\begin{itemize}
    \item Tecnologie web moderne.
    \item Game Development.
    \item Futsal.
  \end{itemize}}
  {Christian lavora su più progetti contemporaneamente e apprezza strumenti integrati e reattivi.
  Usa BugBoard26 per gestire bug durante i test interni. La sua priorità è ridurre al minimo lo sforzo per gestire gli errori,
  un'alta qualità del progetto. Aspetta il suo momento di gloria...}
\end{personaenv}

\begin{personaenv}
  {Marzio Masciarelli}
  {36}
  {Firenze, Italia}
  {sposata}
  {Project Manager}
  {Amministratore}
  {organizzato, determinato, comunicativo}
  {\begin{itemize}
    \item Assegnare bug agli sviluppatori in base al carico.
    \item Monitorare il progresso e le scadenze.
    \item Produrre report mensili per la direzione.
  \end{itemize}}
  {\begin{itemize}
    \item Arte e scultura.
    \item Analisi dei dati e metriche di produttività.
    \item Robotica.
  \end{itemize}}
  {Marzio supervisiona più progetti contemporaneamente. Usa la dashboard di BugBoard26 per monitorare
  le attività e generare report automatici. Predilige strumenti chiari e sintetici per prendere
  decisioni rapide.}
\end{personaenv}

\begin{personaenv}
  {Walter Severino Torrone}
  {44}
  {Napoli, Italia}
  {sposato}
  {Lead Developer}
  {Amministratore}
  {razionale, tecnico, autorevole}
  {\begin{itemize}
    \item Mantenere ordine nella gestione dei bug.
    \item Verificare la qualità delle issue segnalate.
    \item Comunicazione trasparente nel team.
  \end{itemize}}
  {\begin{itemize}
    \item Sistemi informativi multimediali.
    \item Version control e code review.
    \item Software testing e code coverage.
  \end{itemize}}  
  {Andrea gestisce un team di 10 sviluppatori. Usa BugBoard26 per assegnare issue e verificare
  la qualità delle soluzioni. Apprezza l’integrazione con strumenti di CI/CD e la cronologia
  dettagliata delle modifiche.}
\end{personaenv}

\begin{personaenv}
    {Laura Greco}
    {45}
    {Verona, Italia}
    {sposata}
    {Responsabile Qualità}
    {Utente esterno}
    {attenta, diplomatica, metodica}
    {\begin{itemize}
        \item Monitorare bug relativi al software.
        \item Validare la qualità del software consegnato.
        \item Tracciabilità dei problemi aperti.
    \end{itemize}}
    {\begin{itemize}
        \item Sicurezza Informatica.
        \item Tennis.
        \item Politica e attualità.
    \end{itemize}}
    {Laura rappresenta il cliente aziendale e utilizza BugBoard26 per verificare che i problemi
    segnalati siano risolti nei tempi previsti. Non ha competenze tecniche avanzate, quindi
    apprezza dashboard semplici e report chiari.}
\end{personaenv}

\begin{personaenv}
    {Pietro Balzano} 
    {60}
    {Napoli, Italia}
    {sposato}
    {Direttore Tecnico}
    {Utente esterno}
    {esigente, pragmatico, orientato ai risultati}
    {\begin{itemize}
        \item Ottenere una visione globale dell’avanzamento del progetto.
        \item Ricevere report periodici sull’attività del team.
        \item Verificare che le criticità siano risolte in modo tempestivo.
    \end{itemize}}
    {\begin{itemize}
        \item Moda.
        \item Astrofisica.
        \item Disegno tecnico.
    \end{itemize}}
    {Pietro accede a BugBoard26 solo per consultare report e statistiche di progetto.
    Richiede dati sintetici, esportabili e affidabili. Vuole trasparenza e controllo
    sullo stato dei lavori senza entrare nei dettagli tecnici.} 
\end{personaenv}